\documentclass[a4paper]{article}
\usepackage{amsfonts}
\usepackage{amsmath}
\usepackage{amssymb}
\usepackage{graphicx}  
\usepackage{cancel}

\newcommand{\degree}{^{\circ}}
\newcommand{\cis}{\operatorname{cis}}
\newcommand{\Bgsin}{\operatorname{Bgsin}}
\newcommand{\Bgcos}{\operatorname{Bgcos}}
\newcommand{\Bgtan}{\operatorname{Bgtan}}
\newcommand{\limit}{\lim\limits}

\begin{document}
  
\noindent \large Auteur: Rune Suy \\
\noindent \large Studierichting: Informatica\\
\noindent \large Oefeningenreeks 3F nr. 1.3\\

\medskip

\normalsize

$f(x)=\dfrac{x^3 + 2x^2 + 2x - 2}{x^2 + 1}$\\

$x^2 + 1 = 0 \Leftrightarrow x^2 = -1 \Rightarrow$ noemer heeft geen reële nulpunten.\\

Geen VA.\\

$\limit_{x \rightarrow \infty} \dfrac{f(x)}{x} = \limit_{x \rightarrow \infty}\dfrac{x^3 + 2x^2 + 2x - 2}{x^3 + x} = 1 = a$\\

$\limit_{x\rightarrow \infty} (f(x) - ax) = \limit_{x\rightarrow \infty} (f(x) - 1x) = \dfrac{x^3 + 2x^2 + 2x - 2 - (x^3 + x)}{x^2 + 1}$\\

$=\limit_{x\rightarrow \infty} \dfrac{2x^3 + 2x - 2}{x^2 + 1} = \limit_{x\rightarrow \infty}\dfrac{2x^2}{x^2} = 2 = b$\\

SA: $y = x + 2$ bij zowel $+\infty$ als $-\infty$\\

$f(x) - x - 2 = \dfrac{\cancel{x^3} + \cancel{2x^2} + 2x - 2 - (\cancel{x^3} + x) - \cancel{2x^2} - 2}{x^2 + 1} = \dfrac{x-4}{x^2 + 1}$\\

\begin{tabular}{c|c c c}
    $x$         &       & $4$ & \\ \hline
    $f(x)-(x+2)$& $-$   & $0$ & $+$ \\
\end{tabular}

\bigskip

Antwoord: \\

geen HA\\

geen VA\\

wel SA: $y = x + 2$ en $f$ ligt boven SA bij $+\infty$ en onder SA bij $-\infty$\\

\begin{thebibliography}{999}
%\bibitem{a} Referentie
\end{thebibliography}
\end{document} 
